\usepackage{xcolor}
\usepackage{amssymb,amsmath,amsfonts,amsthm,setspace,graphicx,mathtools,quiver,tikz-cd,tikz,ccicons,mathrsfs,adjustbox}
% nfold is loaded here rather than left to quiver, which loads it too. That is
% not redundancy: quiver only grew the \usetikzlibrary{nfold} line recently
% (quiver 2026/07/27 has it), and the TeX Live snapshot in
% ghcr.io/xu-cheng/texlive-alpine:latest still ships a release that does not.
% So ch04.tex's `nfold' arrow key resolved locally and failed in CI with
% "I do not know the key '/tikz/nfold'" — pointing at the chapter, four files
% away from the missing load, with nothing naming quiver at all.
%
% Depending on it here makes the requirement ours and version-independent. The
% failure also changes shape for the better: a missing tikz-nfold now halts in
% the preamble naming the library, instead of surfacing wherever an arrow
% happens to use it.
\usetikzlibrary{nfold}
% keytheorems, not thmtools: thmtools 0.76 builds shared theorem counters with
% its aliasctr package, which LaTeX 2026 broke — `sibling=` then fails with
% "Command \c@proposition already defined". Fix is unreleased upstream
% (muzimuzhi/thmtools#75). thmtools-compat keeps \declaretheorem unchanged;
% loading both packages is a fatal error, so thmtools is gone from the line above.
\usepackage[thmtools-compat]{keytheorems}
\usepackage[hyphens]{url}
\usepackage{hyperref}
% cleveref after hyperref, as it requires.
\usepackage{cleveref}

\definecolor{myblue}{HTML}{0000EE}
\hypersetup{
    %colorlinks=true,
    %linkcolor=myblue,
    hidelinks
    }
% \DocSlug is the topic's folder name, taken from the last segment of \TexRepo
% — the one place a main.tex writes the directory it lives in. It is the PDF
% metadata title, so a viewer's window title and a file manager's Title column
% read "algebraic_k_theory", matching pdf/<topic>.pdf, rather than the typeset
% \DocTitle. \str_set:Ne makes every character catcode 12, which is what keeps
% the underscore in a name like galois_theory from being read as a subscript.
%
% \TexRepo, like \DocTitle, is defined by main.tex *after* this file is
% \input, so both the split and the metadata have to wait until it exists.
% pdflang is deliberately absent: hyperref's luatex driver does not emit /Lang,
% and the \DocumentMetadata route that would is fatal here — it sends
% keytheorems into an infinite loop at \declaretheorem.
\ExplSyntaxOn
\tl_new:N \g_mathstudy_slug_tl
\newcommand{\DocSlug}{\g_mathstudy_slug_tl}
\AtBeginDocument{
    \seq_set_split:Nne \l_tmpa_seq { / } { \TexRepo }
    \str_set:Ne \l_tmpa_str { \seq_item:Nn \l_tmpa_seq { -1 } }
    \tl_gset_eq:NN \g_mathstudy_slug_tl \l_tmpa_str
    \hypersetup{
        pdftitle = { \DocSlug },
        pdfauthor = { \DocAuthorName }
        }
    }
\ExplSyntaxOff

%%% ---- Index ----
% Every defined term in tex/*/ch0N.tex is written with one of these three, and
% nothing else in a chapter is bold: \textbf is the term-defining markup here,
% so the macro that typesets it is also the one that indexes it. The rules —
% which macro, what a reading looks like — are docs/index-convention.md;
% scripts/test_index_markup.py is what notices a chapter that forgot.
%
% All three typeset exactly the bytes the \textbf they replaced did. That is
% the point: converting a chapter changes the index and nothing else, and the
% check is that the PDF is unchanged apart from the new 索引 page.
%
% The reading is mandatory rather than optional because of how its absence
% fails. upmendex sorts a kanji headword with no reading by code point and
% files it in a clump at the end of the index, silently, with no warning in
% the .ilg — so a forgotten reading is invisible in the log and visible only
% to whoever reads the finished index looking for a term that is not where it
% belongs. As a required argument it is a compile error on the line you are
% writing instead.
\usepackage{makeidx}
\makeindex

% \term{<よみ>}{<和文>}{<欧文>} — the ordinary case, and almost every term.
% The optional argument is the sort key for the *English* entry, needed only
% when the gloss carries mathematics: $\lambda$-term has to sort under L, and
% upmendex would otherwise file it under the literal "$".
\NewDocumentCommand{\term}{o m m m}{%
  \textbf{#3}(#4)%
  \index{#2@#3}%
  \IfNoValueTF{#1}{\index{#4}}{\index{#1@#4}}%
}

% \termja{<よみ>}{<和文>} — a term the notes give no English gloss (対象, 射).
% Separate from \term with an empty fourth argument on purpose: "this term has
% no English" and "I forgot the English" must not look alike in the source.
\NewDocumentCommand{\termja}{m m}{%
  \textbf{#2}\index{#1@#2}%
}

% \termen{<欧文>} — a term the notes name only in Latin script (Nakayama's
% lemma, Morita invariance). No reading: ICU sorts Latin without help.
\NewDocumentCommand{\termen}{m}{%
  \textbf{#1}\index{#1}%
}

\crefname{definition}{定義}{定義}
\crefname{proposition}{命題}{命題}
\crefname{lemma}{補題}{補題}
\crefname{theorem}{定理}{定理}
\crefname{corollary}{系}{系}
\crefname{remark}{注意}{注意}
\crefname{example}{例}{例}

\newtheoremstyle{noperiod}
  {}
  {}
  {\normalfont}
  {}
  {\bfseries}
  { }
  {.5em}
  {}

% The three spacing keys are thmtools' defaults, restated because keytheorems
% has different ones; without them every topic repaginates.
\declaretheoremstyle[
  headfont=\normalfont\bfseries,
  bodyfont=\normalfont,
  headpunct={},
  postheadspace={ },
  spaceabove=3pt,
  spacebelow=3pt,
  qed=$\square$
]{style_with_box}

\declaretheorem[name=定義, style=style_with_box, within=section]{definition}
\declaretheorem[name=命題, style=style_with_box, sibling=definition]{proposition}
\declaretheorem[name=補題, style=style_with_box, sibling=definition]{lemma}
\declaretheorem[name=定理, style=style_with_box, sibling=definition]{theorem}
\declaretheorem[name=系,   style=style_with_box, sibling=definition]{corollary}
\declaretheorem[name=注意, style=style_with_box, sibling=definition]{remark}
\declaretheorem[name=例, style=style_with_box, sibling=definition]{example}

\makeatletter
\renewenvironment{proof}[1][\proofname]{\par
  \pushQED{\qed}%
  \normalfont \topsep6\p@\@plus6\p@\relax
  \trivlist
  \item[\hskip\labelsep
        \itshape
    #1 ]\ignorespaces
}{
  \popQED\endtrivlist\@endpefalse
}
\makeatother

%%% ---- Algebra: maps and functors ----
\DeclareMathOperator{\Hom}{Hom}
\DeclareMathOperator{\End}{End}
\DeclareMathOperator{\Aut}{Aut}
\DeclareMathOperator{\Ext}{Ext}
\DeclareMathOperator{\Tor}{Tor}
\DeclareMathOperator{\Ann}{Ann}
\DeclareMathOperator{\Der}{Der}
\DeclareMathOperator{\Frac}{Frac}
\DeclareMathOperator{\Ind}{Ind}
\DeclareMathOperator{\Res}{Res}       % restriction (induction's partner)
\DeclareMathOperator{\res}{res}       % residue, if you want lowercase too

%%% ---- Kernels, images, cokernels ----
\renewcommand{\ker}{\operatorname{Ker}}  % \ker is predefined -> renew
\DeclareMathOperator{\coker}{Coker}
\DeclareMathOperator{\im}{Im}          % NOT \Im: that's Fraktur ℑ
\DeclareMathOperator{\coim}{Coim}

%%% ---- Linear algebra ----
\DeclareMathOperator{\rank}{rank}
\DeclareMathOperator{\trace}{tr}
\DeclareMathOperator{\diag}{diag}
\DeclareMathOperator{\id}{id}
\DeclareMathOperator{\sgn}{sgn}
\DeclareMathOperator{\adj}{adj}
\DeclareMathOperator{\proj}{proj}
\DeclareMathOperator{\spann}{span}        % \span is a TeX primitive
\DeclareMathOperator{\Mat}{Mat}

%%% ---- Matrix / Lie groups ----
\DeclareMathOperator{\GL}{GL}
\DeclareMathOperator{\SL}{SL}
\DeclareMathOperator{\PGL}{PGL}
\DeclareMathOperator{\PSL}{PSL}
\DeclareMathOperator{\OO}{O}              % \O is Ø in text mode
\DeclareMathOperator{\SO}{SO}
\DeclareMathOperator{\UU}{U}
\DeclareMathOperator{\SU}{SU}
\DeclareMathOperator{\Sp}{Sp}
\DeclareMathOperator{\Aff}{Aff}
\DeclareMathOperator{\ad}{ad}
\DeclareMathOperator{\Ad}{Ad}
\DeclareMathOperator{\Lie}{Lie}

%%% ---- Group theory ----
\DeclareMathOperator{\Sym}{Sym}
\DeclareMathOperator{\Alt}{Alt}
\DeclareMathOperator{\Stab}{Stab}
\DeclareMathOperator{\Orb}{Orb}
\DeclareMathOperator{\Inn}{Inn}
\DeclareMathOperator{\Out}{Out}
\DeclareMathOperator{\Gal}{Gal}
\DeclareMathOperator{\ord}{ord}
\DeclareMathOperator{\charac}{char}       % \char is a TeX primitive
\DeclareMathOperator{\Cent}{Z}
\DeclareMathOperator{\Norm}{N}

%%% ---- Category theory ----
\DeclareMathOperator{\Ob}{Ob}
\DeclareMathOperator{\Mor}{Mor}
\DeclareMathOperator{\dom}{dom}
\DeclareMathOperator{\cod}{cod}
\DeclareMathOperator{\Fun}{Fun}
\DeclareMathOperator{\Map}{Map}
\DeclareMathOperator{\op}{op}
\DeclareMathOperator{\ev}{ev}
\DeclareMathOperator{\Nat}{Nat}
\DeclareMathOperator{\Cone}{Cone}
\DeclareMathOperator*{\colim}{colim}      % star: limits go under in display
\DeclareMathOperator*{\hocolim}{hocolim}
\DeclareMathOperator*{\holim}{holim}

%%% ---- Algebraic geometry / commutative algebra ----
\DeclareMathOperator{\Spec}{Spec}
\DeclareMathOperator{\Proj}{Proj}
\DeclareMathOperator{\Supp}{Supp}
\DeclareMathOperator{\supp}{supp}
\DeclareMathOperator{\Div}{Div}           % \div is ÷
\DeclareMathOperator{\Pic}{Pic}
\DeclareMathOperator{\Jac}{Jac}
\DeclareMathOperator{\codim}{codim}
\DeclareMathOperator{\hgt}{ht}     % height of a prime ideal
\DeclareMathOperator{\depth}{depth}

%%% ---- Topology / geometry ----
\DeclareMathOperator{\interior}{int}      % \int is the integral sign
\DeclareMathOperator{\cl}{cl}
\DeclareMathOperator{\bd}{bd}
\DeclareMathOperator{\diam}{diam}
\DeclareMathOperator{\dist}{dist}
\DeclareMathOperator{\vol}{vol}
\DeclareMathOperator{\area}{area}

%%% ---- Analysis / vector calculus ----
\DeclareMathOperator{\grad}{grad}
\DeclareMathOperator{\curl}{curl}
\DeclareMathOperator{\divergence}{div}
\DeclareMathOperator*{\esssup}{ess\,sup}
\DeclareMathOperator*{\essinf}{ess\,inf}
\DeclareMathOperator*{\argmax}{arg\,max}
\DeclareMathOperator*{\argmin}{arg\,min}

%%% ---- Number theory / discrete ----
\DeclareMathOperator{\lcm}{lcm}
\DeclareMathOperator{\Tr}{Tr}             % field trace
\DeclareMathOperator{\Nm}{N}              % field norm
\DeclareMathOperator{\Frob}{Frob}
\DeclareMathOperator{\Li}{Li}

%%% ---- Probability / statistics ----
\DeclareMathOperator{\E}{\mathbb{E}}
\DeclareMathOperator{\Var}{Var}
\DeclareMathOperator{\Cov}{Cov}
\DeclareMathOperator{\Corr}{Corr}
\DeclareMathOperator{\Bias}{Bias}

\renewcommand{\proofname}{\bfseries 証明}
\renewcommand{\qedsymbol}{$\blacksquare$}

\newcommand{\GitHub}{https://github.com/ys-math}
% \DocAuthorName is the bare text; \DocAuthor is the same name wrapped in a
% link to \GitHub, and is what \author{} and the colophon print. The two are
% separate because the PDF's /Author metadata (\hypersetup above) must stay
% plain text — an \href there would be typeset into the metadata string.
\newcommand{\DocAuthorName}{ys-math}
\newcommand{\DocAuthor}{\href{\GitHub}{\DocAuthorName}}
\newcommand{\DocRepo}{https://github.com/ys-math/math-study.git}
\newcommand{\DocCompiler}{LuaLaTeX}
\newcommand{\DocClass}{jlreq}